% Part: applied-modal-logics
% Chapter: temporal-logic
% Section: temporal-logic-semantics

\documentclass[../../../include/open-logic-section]{subfiles}

\begin{document}

\olfileid{aml}{tl}{sem}

\olsection{دلالة المنطق الزمني}

\begin{defn}
تتركب اللغة الأساسية للمنطق الزمني مما يأتي:
\begin{enumerate}
  \tagitem{prvFalse}{الثابت القضي الدال على !!{falsity}~$\lfalse$.}{}
  \tagitem{prvTrue}{الثابت القضي الدال على !!{truth}~$\ltrue$.}{}
  \item مجموعة !!{denumerable} من ال!!{propositional variable}s:
    $\Obj p_{\OLMathNumeral{0}}$، $\Obj p_{\OLMathNumeral{1}}$، $\Obj p_{\OLMathNumeral{2}}$، \dots
  \item الروابط القضية: \startycommalist
  \iftag{prvNot}{\ycomma $\lnot$ (النفي)}{}%
  \iftag{prvAnd}{\ycomma $\land$ (العطف)}{}%
  \iftag{prvOr}{\ycomma $\lor$ (الفصل)}{}%
  \iftag{prvIf}{\ycomma $\lif$ (!!{conditional})}{}%
  \iftag{prvIff}{\ycomma $\liff$ (!!{biconditional})}{}.
  \item مؤثرا الماضي~$\Ptemp$ و$\Htemp$.
  \item مؤثرا المستقبل~$\Ftemp$ و$\Gtemp$.
\end{enumerate}
\end{defn}

وسيأتي بعد هذا ذكر ما يجوز أن يزاد عليها من ضروب المؤثرات الجهوية.

\begin{defn}
وأما \emph{!!^{formula}s} اللغة الزمنية فتحدّ استقرائيًا على الوجه الآتي:
\begin{enumerate}
\tagitem{prvFalse}{$\lfalse$ !!{formula} ذرية.}{}

\tagitem{prvTrue}{$\ltrue$ !!{formula} ذرية.}{}

\item كل متغير قضي~$\Obj p_{\OLId{latin-ordinary}{i}}$ هو !!{formula} (ذرية).

\tagitem{prvNot}{إذا كانت~$!A$ !!a{formula}، كانت~$\lnot !A$
  !!a{formula}.}{}

\tagitem{prvAnd}{إذا كانت~$!A$ و$!B$ !!{formula}s، كانت
  $(!A \land !B)$ !!a{formula}.}{}

\tagitem{prvOr}{إذا كانت~$!A$ و$!B$ !!{formula}s، كانت
  $(!A \lor !B)$ !!a{formula}.}{}

\tagitem{prvIf}{إذا كانت~$!A$ و$!B$ !!{formula}s، كانت
  $(!A \lif !B)$ !!a{formula}.}{}

\tagitem{prvIff}{إذا كانت~$!A$ و$!B$ !!{formula}s، كانت
  $(!A \liff !B)$ !!a{formula}.}{}

% تصحيح مصدر (0478-I1): استُعمل الماكرو الزمني Ftemp بدل الحرف F العاري الوارد مرة في الأصل الإنجليزي.
\item إذا كانت~$!A$ !!a{formula}، فإن~$\Ptemp !A$ و$\Htemp !A$
  و$\Ftemp !A$ و$\Gtemp !A$ كلها !!{formula}s.

\tagitem{limitClause}{لا شيء غير ذلك !!a{formula}.}{}
\end{enumerate}
\end{defn}

وتفسر هذه اللغة بنماذج علائقية، على نحو ما تصنع سائر ضروب المنطق المقصدي.

\begin{defn}
  \emph{نموذج} اللغة الزمنية هو الثلاثي
  $\mModel{M}=\tuple{{\OLId{latin-looped}{T}},\prec,{\OLId{latin-looped}{V}}}$، وفيه:
  \begin{enumerate}
  \item ${\OLId{latin-looped}{T}}$ مجموعة غير فارغة، تُفسر عناصرها بوصفها نقاطًا في الزمن.
  \item $\prec$ علاقة ثنائية على~${\OLId{latin-looped}{T}}$.
  \item ${\OLId{latin-looped}{V}}$ دالة تسند إلى كل~!!{propositional
    variable}~${\OLId{latin-ordinary}{p}}$ مجموعة~${\OLId{latin-looped}{V}}({\OLId{latin-ordinary}{p}})$ من النقاط الزمنية.
  \end{enumerate}
  فإذا كان~${\OLId{latin-ordinary}{t}} \prec {\OLId{latin-ordinary}{t}}'$ قلنا إن~${\OLId{latin-ordinary}{t}}$ \emph{يسبق}~${\OLId{latin-ordinary}{t}}'$؛ وإذا كان
  ${\OLId{latin-ordinary}{t}} \in {\OLId{latin-looped}{V}}({\OLId{latin-ordinary}{p}})$ قلنا إن~${\OLId{latin-ordinary}{p}}$ \emph{صادق عند}~${\OLId{latin-ordinary}{t}}$.
\end{defn}

ولا نفرض في هذا الموضع شرطًا على علاقة السبق~$\prec$؛ فلذلك لا تتقيد بنية
النموذج الزمني عندنا الآن: فقد يكون الزمن فيه خطيًا، أو متفرعًا، أو دائريًا،
أو على أي هيئة أخرى.

وكما في المنطق الجهوي النظامي، يعين كل نموذج زمني ال!!{formula}s الصادقة
عند كل نقطة منه. ونعبّر عن أصل هذه الدلالة العلائقية بقولنا: «يصدق
النموذج~$\mModel{M}$ !!{formula}~$!A$ عند النقطة~${\OLId{latin-ordinary}{t}}$». وتحد علاقة الصدق
استقرائيًا؛ وهي في المؤثرات غير الجهوية عين العلاقة المتقدمة في المنطق
الجهوي النظامي.

\begin{defn}\ollabel{defn:tmodels}
  \emph{صدق !!a{formula}~$!A$ عند~${\OLId{latin-ordinary}{t}}$} في~$\mModel M$، ورمزه
  $\mSat{{\OLId{latin-looped}{M}}}{!A}[{\OLId{latin-ordinary}{t}}]$، يحد استقرائيًا بما يأتي:
  \begin{enumerate}
  \tagitem{prvFalse}{\indcase{!A}{\lfalse}{لا يتحقق أبدًا
    $\mSat{{\OLId{latin-looped}{M}}}{\lfalse}[{\OLId{latin-ordinary}{t}}]$}.}{}
  \tagitem{prvTrue}{\indcase{!A}{\ltrue}{يتحقق دائمًا
    $\mSat{{\OLId{latin-looped}{M}}}{\ltrue}[{\OLId{latin-ordinary}{t}}]$}.}{}
  \item $\mSat{{\OLId{latin-looped}{M}}}{{\OLId{latin-ordinary}{p}}}[{\OLId{latin-ordinary}{t}}]$ إذا وفقط إذا كان~${\OLId{latin-ordinary}{t}} \in {\OLId{latin-looped}{V}}({\OLId{latin-ordinary}{p}})$.
  \tagitem{prvNot}{\indcase{!A}{\lnot !B}{$\mSat{{\OLId{latin-looped}{M}}}{\indfrm}[{\OLId{latin-ordinary}{t}}]$ إذا
    وفقط إذا كان~$\mSat/{{\OLId{latin-looped}{M}}}{!B}[{\OLId{latin-ordinary}{t}}]$}.}{}
  \tagitem{prvAnd}{\indcase{!A}{(!B \land !C)}{$\mSat{{\OLId{latin-looped}{M}}}{\indfrm}[{\OLId{latin-ordinary}{t}}]$
    إذا وفقط إذا كان~$\mSat{{\OLId{latin-looped}{M}}}{!B}[{\OLId{latin-ordinary}{t}}]$ و$\mSat{{\OLId{latin-looped}{M}}}{!C}[{\OLId{latin-ordinary}{t}}]$}.}{}
  \tagitem{prvOr}{\indcase{!A}{(!B \lor !C)}{$\mSat{{\OLId{latin-looped}{M}}}{\indfrm}[{\OLId{latin-ordinary}{t}}]$
    إذا وفقط إذا كان~$\mSat{{\OLId{latin-looped}{M}}}{!B}[{\OLId{latin-ordinary}{t}}]$ أو~$\mSat{{\OLId{latin-looped}{M}}}{!C}[{\OLId{latin-ordinary}{t}}]$}
    (أو كلاهما).}{}
  \tagitem{prvIf}{\indcase{!A}{(!B \lif !C)}{$\mSat{{\OLId{latin-looped}{M}}}{\indfrm}[{\OLId{latin-ordinary}{t}}]$
    إذا وفقط إذا كان~$\mSat/{{\OLId{latin-looped}{M}}}{!B}[{\OLId{latin-ordinary}{t}}]$ أو~$\mSat{{\OLId{latin-looped}{M}}}{!C}[{\OLId{latin-ordinary}{t}}]$}.}{}
  \tagitem{prvIff}{\indcase{!A}{(!B \liff !C)}{$\mSat{{\OLId{latin-looped}{M}}}{\indfrm}[{\OLId{latin-ordinary}{t}}]$
    إذا وفقط إذا تحقق كل من~$\mSat{{\OLId{latin-looped}{M}}}{!B}[{\OLId{latin-ordinary}{t}}]$ و$\mSat{{\OLId{latin-looped}{M}}}{!C}[{\OLId{latin-ordinary}{t}}]$، أو
    لم يتحقق أي من~$\mSat{{\OLId{latin-looped}{M}}}{!B}[{\OLId{latin-ordinary}{t}}]$ و$\mSat{{\OLId{latin-looped}{M}}}{!C}[{\OLId{latin-ordinary}{t}}]$}.}{}
  \item\ollabel{defn:sub:mmodels-p}
    \indcase{!A}{\Ptemp !B}{$\mSat{{\OLId{latin-looped}{M}}}{\indfrm}[{\OLId{latin-ordinary}{t}}]$ إذا وفقط إذا كان
    $\mSat{{\OLId{latin-looped}{M}}}{!B}[{\OLId{latin-ordinary}{t}}']$ لبعض~${\OLId{latin-ordinary}{t}}' \in {\OLId{latin-looped}{T}}$ بحيث~${\OLId{latin-ordinary}{t}}' \prec {\OLId{latin-ordinary}{t}}$}
  \item\ollabel{defn:sub:mmodels-h}
    \indcase{!A}{\Htemp !B}{$\mSat{{\OLId{latin-looped}{M}}}{\indfrm}[{\OLId{latin-ordinary}{t}}]$ إذا وفقط إذا كان
    $\mSat{{\OLId{latin-looped}{M}}}{!B}[{\OLId{latin-ordinary}{t}}']$ لكل~${\OLId{latin-ordinary}{t}}' \in {\OLId{latin-looped}{T}}$ بحيث~${\OLId{latin-ordinary}{t}}' \prec {\OLId{latin-ordinary}{t}}$}
  \item\ollabel{defn:sub:mmodels-f}
    \indcase{!A}{\Ftemp !B}{$\mSat{{\OLId{latin-looped}{M}}}{\indfrm}[{\OLId{latin-ordinary}{t}}]$ إذا وفقط إذا كان
    $\mSat{{\OLId{latin-looped}{M}}}{!B}[{\OLId{latin-ordinary}{t}}']$ لبعض~${\OLId{latin-ordinary}{t}}' \in {\OLId{latin-looped}{T}}$ بحيث~${\OLId{latin-ordinary}{t}} \prec {\OLId{latin-ordinary}{t}}'$}
  \item\ollabel{defn:sub:mmodels-g}
    \indcase{!A}{\Gtemp !B}{$\mSat{{\OLId{latin-looped}{M}}}{\indfrm}[{\OLId{latin-ordinary}{t}}]$ إذا وفقط إذا كان
    $\mSat{{\OLId{latin-looped}{M}}}{!B}[{\OLId{latin-ordinary}{t}}']$ لكل~${\OLId{latin-ordinary}{t}}' \in {\OLId{latin-looped}{T}}$ بحيث~${\OLId{latin-ordinary}{t}} \prec {\OLId{latin-ordinary}{t}}'$}
  \end{enumerate}
\end{defn}

ومن هذه الدلالة يظهر أن~$\Ptemp$ و$\Htemp$ متقابلان، وأن~$\Ftemp$
و$\Gtemp$ كذلك. ولذلك يصح أن نعرّف~$\Htemp !A$ بأنه
$\lnot \Ptemp \lnot !A$، وأن نعرّف~$\Gtemp$ من~$\Ftemp$ على نظيره.

\end{document}
